\section{Results}
\label{sec:results}

We present results in the order established by our research questions: existence (RQ1), architecture
ranking (RQ2), gap-specific signal (RQ3), and differential advantage (RQ4). The evidence builds toward
our central finding: generalization gap is a learnable, architecturally-sensitive, and genuinely distinct
weight-space prediction target.

\subsection{RQ1: Gap Learnability and Target Independence}

Figure~\ref{fig:gap_spearman} shows Spearman rank correlations for all four encoders on the gap prediction
target. Two encoders exceed the existence threshold ($r > 0.5$): FlatMLP ($r = 0.5567$) and DWSNet
($r = 0.5104$). This establishes that generalization gap is a learnable signal in weight tensors --- accessible
even to a position-indexed flat encoder without equivariant constraints.

\begin{table}[h]
\centering
\caption{Gap and Test Accuracy Spearman (Test Split, $N=1{,}000$).}
\label{tab:results}
\begin{tabular}{@{}lllll@{}}
\toprule
Encoder & Spearman(gap) & 95\% CI & Spearman(test\_acc) & 95\% CI \\
\midrule
FlatMLP & 0.5567$^\dagger$ & [0.4850, 0.5801]$^\ddagger$ & 0.2790 & [0.2173, 0.3343] \\
DWSNet  & 0.5104$^\dagger$ & --- & 0.4553 & [0.4020, 0.5054] \\
NFT     & 0.5752 & [0.5339, 0.6158] & 0.4801 & [0.4326, 0.5262] \\
GNN     & 0.3747 & [0.3180, 0.4265] & 0.3480 & [0.2913, 0.4037] \\
\bottomrule
\end{tabular}
\end{table}

\noindent\textit{$^\dagger$ FlatMLP and DWSNet gap point estimates in Table~\ref{tab:results} are from the existence run (h-e1). FlatMLP gap in the architecture comparison run (h-m1) is 0.5330; this h-m1 value is used as the control baseline in the $\Delta$ computation (Table~\ref{tab:delta}). $^\ddagger$ FlatMLP 95\% CI [0.4850, 0.5801] from h-m1 bootstrap.}

The A1 audit confirms that gap is not trivially derivable from test accuracy: Spearman(gap, $-$test\_acc)
$= -0.142$, far below the collinearity threshold of 0.95. Predicting gap is a genuinely independent task.

\textbf{The counterintuitive asymmetry.} Figure~\ref{fig:dual_target} reveals that FlatMLP
predicts gap ($r = 0.557$) almost twice as well as test accuracy ($r = 0.279$) from identical weight inputs.
This reverses the expectation from prior literature (where FlatMLP achieves $r \approx 0.85$ on test\_acc in the
original Unterthiner zoo). We discuss this anomaly in Section~\ref{sec:discussion}; crucially, it does not affect the
gap prediction results, which are our primary contribution.

\subsection{RQ2: Architecture Ranking on Gap Prediction}

\textbf{NFT achieves the highest gap Spearman} ($r = 0.5752$) with a 95\% CI of [0.5339, 0.6158]. For
comparison, FlatMLP's 95\% CI from the same run (h-m1) is [0.4850, 0.5801]. The two intervals overlap
at the margin, but NFT's lower bound (0.5339) exceeds FlatMLP's point estimate (0.5330), providing
directional evidence of NFT's advantage. Figure~\ref{fig:gap_spearman} shows this ranking clearly: NFT $>$ FlatMLP $>$
DWSNet $>$ GNN for gap prediction.

The architecture-specificity of this result is the key finding. DWSNet (within-layer equivariance)
achieves $r = 0.4881$ on gap --- below FlatMLP's $r = 0.5330$. GNN achieves only $r = 0.3747$, the lowest
among all encoders. This means equivariance is not uniformly helpful for gap prediction: only the
cross-layer attention architecture (NFT) provides a meaningful advantage, while within-layer and
graph-structured equivariance actually underperform the non-equivariant baseline.

\subsection{RQ3: Gap-Specific Signal Independent of Test Accuracy (P3)}

This is the paper's central result. The partial Spearman analysis tests whether NFT's gap predictions
contain information about true gap that cannot be explained by test accuracy rank alone.

\textbf{NFT partial Spearman (gap $|$ test\_acc): $r = 0.7305$, $p = 1.60 \times 10^{-167}$}

Figure~\ref{fig:partial_corr} shows the residual scatter for the P3 analysis. After
partialling out test accuracy rank from both NFT's gap predictions and the true gap labels, the
residual correlation remains $r = 0.73$. This result is:

\begin{enumerate}
\item \textbf{Stronger than the direct Spearman} ($0.73 > 0.57$): the gap-specific component of NFT's
   predictions is more strongly correlated with the gap-specific component of true gap than the
   combined signal. Partialling out test accuracy isolates an overfitting structure that NFT captures
   particularly well.

\item \textbf{Statistically unambiguous} ($p = 1.60 \times 10^{-167}$): with $N = 1{,}000$ test samples, the p-value is
   not a borderline result --- the gap-specific signal is real and large.

\item \textbf{Architecture-specific}: this analysis was performed for NFT as the encoder with the highest
   direct gap Spearman. The result confirms that NFT's cross-layer attention is not merely learning
   a test-accuracy proxy --- it captures genuine overfitting structure in weight space.
\end{enumerate}

This result answers RQ3: yes, gap predictions from NFT contain substantial information about true gap
that is independent of test accuracy. Generalization gap encodes a distinct signal in weight space that
cross-layer attention is well-positioned to extract.

\subsection{RQ4: Differential Advantage ($\Delta$) --- Null Result}

Figure~\ref{fig:delta} shows the $\Delta$ values for equivariant encoders with 95\% bootstrap CIs.

\begin{table}[h]
\centering
\caption{Differential Advantage $\Delta = [\text{gap improvement}] - [\text{test\_acc improvement}]$ over FlatMLP.}
\label{tab:delta}
\begin{tabular}{@{}llllll@{}}
\toprule
Encoder & gap Spearman & acc Spearman & $\Delta$ & 95\% CI & $\Delta > 0.02$? \\
\midrule
DWSNet & 0.4881 & 0.4553 & $-0.2212$ & entirely $< 0$ & No \\
NFT    & 0.5752 & 0.4801 & $-0.1589$ & entirely $< 0$ & No \\
GNN    & 0.3747 & 0.3480 & $-0.2272$ & entirely $< 0$ & No \\
\bottomrule
\end{tabular}
\end{table}

\noindent\textit{Gate criterion: $\Delta > 0.02$ for $\geq 2$ encoders. Result: 0/3 encoders pass. Gate: FAIL.}

All equivariant $\Delta$ values are strongly negative, with 95\% CIs entirely below zero. The differential
advantage hypothesis (equivariant encoders improve more on gap than on test\_acc relative to FlatMLP)
is not confirmed under these experimental conditions.

\textbf{Why are $\Delta$ values negative?} The dominant driver is the acc improvement
component --- all equivariant encoders show substantially larger test\_acc improvement over FlatMLP than
gap improvement. Since FlatMLP test\_acc ($r = 0.279$) is anomalously low (vs.\ literature $\approx 0.85$), the
apparent acc improvement for equivariant encoders is inflated. DWSNet shows $0.455 - 0.279 = +0.176$
improvement on test\_acc, while its gap improvement is only $0.488 - 0.533 = -0.045$ (DWSNet is actually
worse than FlatMLP on gap). This confound prevents a clean interpretation of the $\Delta$ result.

\textbf{What we can conclude:} (1) Under our experimental conditions, equivariant encoders do not show a
gap-specific differential advantage. (2) The most likely confounder is the FlatMLP test\_acc baseline
being severely underestimated due to our limited search budget. (3) The gap prediction results (RQ1--RQ3)
are not affected by this issue --- they rely only on gap Spearman, which is internally consistent across
independent runs (FlatMLP: 0.5567 in h-e1, 0.5330 in h-m1; gap $< 5\%$ deviation).
